% sec9_3.tex — Section 9.3: Dickey-Fuller Tests

In this and the following section, we will present several tests of the null hypothesis
that the data generating process (DGP) is I(1). In this section, the I(1) process is a
random walk with or without drift, and the tests are based on the OLS estimation
of appropriate AR(1) equations. Those tests were developed by Dickey (1976) and
Fuller (1976) and are referred to as \textbf{Dickey-Fuller tests} or \textbf{DF tests}. We take the
convention that the sample has $T + 1$ observations $(y_0, y_1, \ldots, y_T)$, so that the
AR(1) equation can be estimated for $t = 1, 2, \ldots, T$.

\subsection*{The AR(1) Model}

The Dickey-Fuller test statistics are derived from the estimation of the first-order
autoregressive model:
\begin{equation}
\label{eq:9.3.1}
y_t = \rho y_{t-1} + \varepsilon_t, \quad \{\varepsilon_t\} \text{ independent white noise.}
\end{equation}
(The case with intercept will be studied in a moment.) This model includes both
I(0) and I(0) processes: if $\rho = 1$, that is, if the associated polynomial equation
$1 - \rho z = 0$ has a unit root, then $\Delta y_t = \varepsilon_t$ and $\{y_t\}$ is a driftless random walk,
whereas if $|\rho| < 1$, the process is zero-mean stationary AR(1). Thus, the hypothesis that the data are generated by a driftless random walk can be formulated as the
null hypothesis that $\rho = 1$. If we take the position that the DGP is either I(0) or
I(1), then $-1 < \rho \leq 1$. Given this \textit{a priori} restriction on $\rho$, the I(0) alternative can
be expressed as $\rho < 1$. So one-tailed tests will be used to test the null of $\rho = 1$
against the I(0) alternative.

The tests will be based on $\hat{\rho}$, the OLS estimate of $\rho$ in the AR(1) equation~\eqref{eq:9.3.1}. Under the I(1) null of $\rho = 1$, the sampling error is $\hat{\rho} - 1$, whose expression
is given by
\begin{equation}
\label{eq:9.3.2}
\hat{\rho} - 1 \equiv \frac{\sum_{t=1}^{T} y_t\, y_{t-1}}{\sum_{t=1}^{T} (y_{t-1})^2} - 1
= \frac{\sum_{t=1}^{T} \Delta y_t\, y_{t-1}}{\sum_{t=1}^{T} (y_{t-1})^2}.
\end{equation}
As will be verified in a moment, the numerator has a limiting distribution if divided
by $T$ and the denominator has one if divided by $T^2$. So consider the sampling error
magnified by $T$:
\begin{equation}
\label{eq:9.3.3}
T \cdot (\hat{\rho} - 1) = T \cdot \frac{\sum_{t=1}^{T} \Delta y_t\, y_{t-1}}{\sum_{t=1}^{T} (y_{t-1})^2}
= \frac{\frac{1}{T}\sum_{t=1}^{T} \Delta y_t\, y_{t-1}}{\frac{1}{T^2}\sum_{t=1}^{T} (y_{t-1})^2}.
\end{equation}

\subsection*{Deriving the Limiting Distribution under the I(1) Null}

Deriving the limiting distribution of $T \cdot (\hat{\rho} - 1)$ under the I(1) null (a driftless random walk) is very straightforward. Since $\{y_t\}$ is driftless I(1), parts~(a) and~(b) of
Proposition~9.2 are applicable with $\xi_t = y_t$. Furthermore, since $\Delta y_t$ is independent
white noise, $\lambda^2$ (the long-run variance of $\{\Delta y_t\}$) equals $\gamma_0$ (the variance of $\Delta y_t$).
Thus,
\begin{align}
w_{1T} &\equiv \frac{1}{T} \sum_{t=1}^{T} \Delta y_t\, y_{t-1}
\;\dto\; \frac{\gamma_0}{2} W(1)^2 - \frac{\gamma_0}{2} \equiv w_1,
\label{eq:9.3.4} \\
w_{2T} &\equiv \frac{1}{T^2} \sum_{t=1}^{T} (y_{t-1})^2
\;\dto\; \gamma_0 \int_0^1 W(r)^2\,\mathrm{d}r \equiv w_2.
\label{eq:9.3.5}
\end{align}
As noted in the Proposition, the convergence is \emph{joint}, so $\mathbf{w}_T \equiv (w_{1T}, w_{2T})'$ converges
to a $2 \times 1$ random vector, $\mathbf{w} \equiv (w_1, w_2)'$. Since~\eqref{eq:9.3.3} is a
continuous function of $\mathbf{w}_T$ (it equals $w_{1T}/w_{2T}$), it follows from Lemma~2.3(b) that
\begin{equation}
\label{eq:9.3.6}
T \cdot (\hat{\rho} - 1)
= \frac{w_{1T}}{w_{2T}}
\;\dto\; \frac{w_1}{w_2}
= \frac{\frac{\gamma_0}{2} W(1)^2 - \frac{\gamma_0}{2}}{\gamma_0 \int_0^1 W(r)^2\,\mathrm{d}r}
= \frac{\tfrac{1}{2}(W(1)^2 - 1)}{\int_0^1 W(r)^2\,\mathrm{d}r}
\equiv DF_\rho.
\end{equation}

The test statistic $T \cdot (\hat{\rho} - 1)$ is called the \textbf{Dickey-Fuller}, or \textbf{DF, $\bm{\rho}$ statistic}. Several
points are worth noting:

\begin{itemize}
\item $T \cdot (\hat{\rho} - 1)$, rather than the usual $\sqrt{T} \cdot (\hat{\rho} - 1)$, has a nondegenerate limiting
distribution. The estimator $\hat{\rho}$ is said to be \textbf{superconsistent} because it converges
to the true value of unity at a faster rate ($T$).

\item The null hypothesis does not specify the values of $\Var(\varepsilon_t)$ and $y_0$ (the initial
value) of the DGP, yet they do not affect the limiting distribution, the distribution
of the random variable $DF_\rho$. That is, the limiting distribution does not involve
those \textbf{nuisance parameters}. So we can use $T \cdot (\hat{\rho} - 1)$ as the test statistic. This
test of a driftless random walk is called the \textbf{Dickey-Fuller $\bm{\rho}$ test}.

\item Note that both the numerator and the denominator involve the \emph{same} Wiener
process, so they are correlated. Since $W(1)$ is distributed as standard normal,
the numerator of the expression for $DF_\rho$ could be written as $(\chi^2(1) - 1)/2$. But
this would obscure the point being made here.
\end{itemize}

The usual $t$-test for the null of $\rho = 1$, too, has a limiting distribution. It is just
a matter of simple algebra to show that the $t$-value for the null of $\rho = 1$ can be
written as
\begin{equation}
\label{eq:9.3.7}
t = \frac{\hat{\rho} - 1}{s \div \sqrt{\sum_{t=1}^{T} (y_{t-1})^2}}
= \frac{\frac{1}{T}\sum_{t=1}^{T} \Delta y_t\, y_{t-1}}{s \cdot \sqrt{\frac{1}{T^2}\sum_{t=1}^{T}(y_{t-1})^2}},
\end{equation}
where $s$ is the standard error of regression:
\begin{equation}
\label{eq:9.3.8}
s \equiv \sqrt{\frac{1}{T-1}\sum_{t=1}^{T}(y_t - \hat{\rho}\, y_{t-1})^2}.
\end{equation}
It is easy to prove (see Review Question~3) that $s^2$ is consistent for $\gamma_0$ ($\equiv \Var(\Delta y_t)$).
It is then immediate from this and~\eqref{eq:9.3.4} and~\eqref{eq:9.3.5} that
\begin{equation}
\label{eq:9.3.9}
t \;\dto\; \frac{\tfrac{1}{2}(W(1)^2 - 1)}{\sqrt{\int_0^1 W(r)^2\,\mathrm{d}r}}
\equiv DF_t.
\end{equation}
This limiting distribution, too, is free from nuisance parameters. So the $t$-value
can be used for testing, but the critical values for the standard normal distribution
cannot be used; they have to come from a tabulation (provided below) of $DF_t$. This
test is called the \textbf{DF $\bm{t}$ test}.

We summarize the discussion so far as

\begin{proposition}[Dickey-Fuller tests of a driftless random walk, the case without intercept]
\label{prop:9.3}
Suppose that $\{y_t\}$ is a driftless random walk (so $\{\Delta y_t\}$ is independent white noise) with $\E(y_0^2) < \infty$. Consider the regression of $y_t$ on $y_{t-1}$
(without intercept) for $t = 1, 2, \ldots, T$. Then
\begin{align}
&\textit{DF $\rho$ statistic:}\quad T \cdot (\hat{\rho} - 1) \;\dto\; DF_\rho, \notag \\
&\textit{DF $t$ statistic:}\quad t \;\dto\; DF_t, \notag
\end{align}
where $\hat{\rho}$ is the OLS estimate of the $y_{t-1}$ coefficient, $t$ is the $t$-value for the hypothesis that the $y_{t-1}$ coefficient is~$1$, and $DF_\rho$ and $DF_t$ are the random variables defined
in~\eqref{eq:9.3.6} and~\eqref{eq:9.3.9}.
\end{proposition}

Thus, we have two test statistics, $T \cdot (\hat{\rho} - 1)$ and the $t$-value, for the same I(1) null.
Comparison of the (generalizations of) these DF tests will be discussed in the next
section.

\begin{table}[htbp]
\centering
\caption{Critical Values for the Dickey-Fuller $\rho$ Test}
\label{tab:9.1}
\small
\begin{tabular}{l*{8}{r}}
\toprule
Sample & \multicolumn{8}{c}{Probability that the statistic is less than entry} \\
\cmidrule(l){2-9}
size ($T$) & 0.01 & 0.025 & 0.05 & 0.10 & 0.90 & 0.95 & 0.975 & 0.99 \\
\midrule
\multicolumn{9}{c}{Panel (a): $T \cdot (\hat{\rho} - 1)$} \\
\midrule
25   & $-11.8$ & $-9.3$ & $-7.3$ & $-5.3$ & 1.01 & 1.41 & 1.78 & 2.28 \\
50   & $-12.8$ & $-9.9$ & $-7.7$ & $-5.5$ & 0.97 & 1.34 & 1.69 & 2.16 \\
100  & $-13.3$ & $-10.2$ & $-7.9$ & $-5.6$ & 0.95 & 1.31 & 1.65 & 2.09 \\
250  & $-13.6$ & $-10.4$ & $-8.0$ & $-5.7$ & 0.94 & 1.29 & 1.62 & 2.05 \\
500  & $-13.7$ & $-10.4$ & $-8.0$ & $-5.7$ & 0.93 & 1.28 & 1.61 & 2.04 \\
$\infty$ & $-13.8$ & $-10.5$ & $-8.1$ & $-5.7$ & 0.93 & 1.28 & 1.60 & 2.03 \\
\midrule
\multicolumn{9}{c}{Panel (b): $T \cdot (\hat{\rho}^\mu - 1)$} \\
\midrule
25   & $-17.2$ & $-14.6$ & $-12.5$ & $-10.2$ & $-0.76$ & 0.00 & 0.65 & 1.39 \\
50   & $-18.9$ & $-15.7$ & $-13.3$ & $-10.7$ & $-0.81$ & $-0.07$ & 0.53 & 1.22 \\
100  & $-19.8$ & $-16.3$ & $-13.7$ & $-11.0$ & $-0.83$ & $-0.11$ & 0.47 & 1.14 \\
250  & $-20.3$ & $-16.7$ & $-13.9$ & $-11.1$ & $-0.84$ & $-0.13$ & 0.44 & 1.08 \\
500  & $-20.5$ & $-16.8$ & $-14.0$ & $-11.2$ & $-0.85$ & $-0.14$ & 0.42 & 1.07 \\
$\infty$ & $-20.7$ & $-16.9$ & $-14.1$ & $-11.3$ & $-0.85$ & $-0.14$ & 0.41 & 1.05 \\
\midrule
\multicolumn{9}{c}{Panel (c): $T \cdot (\hat{\rho}^\tau - 1)$} \\
\midrule
25   & $-22.5$ & $-20.0$ & $-17.9$ & $-15.6$ & $-3.65$ & $-2.51$ & $-1.53$ & $-0.46$ \\
50   & $-25.8$ & $-22.4$ & $-19.7$ & $-16.8$ & $-3.71$ & $-2.60$ & $-1.67$ & $-0.67$ \\
100  & $-27.4$ & $-23.7$ & $-20.6$ & $-17.5$ & $-3.74$ & $-2.63$ & $-1.74$ & $-0.76$ \\
250  & $-28.5$ & $-24.4$ & $-21.3$ & $-17.9$ & $-3.76$ & $-2.65$ & $-1.79$ & $-0.83$ \\
500  & $-28.9$ & $-24.7$ & $-21.5$ & $-18.1$ & $-3.76$ & $-2.66$ & $-1.80$ & $-0.86$ \\
$\infty$ & $-29.4$ & $-25.0$ & $-21.7$ & $-18.3$ & $-3.77$ & $-2.67$ & $-1.81$ & $-0.88$ \\
\bottomrule
\multicolumn{9}{l}{\footnotesize\textsc{Source}: Fuller (1976, Table 8.5.1), corrected in Fuller (1996, Table 10.A.1).}
\end{tabular}
\end{table}

Panel~(a) of Table~9.1 has critical values for the finite-sample distribution of the
random variable $T \cdot (\hat{\rho} - 1)$ for various sample sizes. The solid curve in Figure~9.3
graphs the finite-sample distribution of $\hat{\rho}$ for $T = 100$, which shows that the OLS
estimate of $\rho$ is biased downward (i.e., the mean of $\hat{\rho}$ is less than~1). These are
obtained by Monte Carlo assuming that $\varepsilon_t$ is Gaussian and $y_0 = 0$. By Proposition~9.3, as the sample size $T$ increases, those critical values converge to the critical
values for $DF_\rho$, reported in the row for $T = \infty$. (Since $y_0$ does
not have to be zero and $\varepsilon_t$ does not have to be Gaussian for the asymptotic results
in Proposition~9.3 to hold, we should obtain the same limiting critical values if the
finite-sample critical values were calculated with a different assumption about the
initial value $y_0$ and the distribution of $\varepsilon_t$.) From the critical values for $T = \infty$,
we can see that the $DF_\rho$ distribution, too, is skewed to the left, with the 5 percent
lower tail critical value of $-8.1$ (so $\text{Prob}(DF_\rho < -8.1) = 5\%$) and the 5 percent
upper tail critical value of~1.28. The small-sample bias in $\hat{\rho}$ is reflected in the
skewness in the limiting distribution of $T \cdot (\hat{\rho} - 1)$.

\begin{figure}[htbp]
\centering
\framebox[\textwidth]{\rule{0pt}{2.5in}\parbox{0.8\textwidth}{\centering
  Figure 9.3: Finite-Sample Distribution of OLS Estimate\\
  of AR(1) Coefficient, $T = 100$\\[6pt]
  --- without intercept or time \quad $\cdots$ with intercept \quad - - - with time}}
\caption{Finite-Sample Distribution of OLS Estimate of AR(1) Coefficient, $T = 100$}
\label{fig:9.3}
\end{figure}

Critical values for $DF_t$ (the limiting distribution of the $t$-value) are in the row
for $T = \infty$ of Table~9.2, Panel~(a). Like $DF_\rho$, it is skewed to the left. The table
also has critical values for finite-sample sizes. Again, they assume that $y_0 = 0$ and
$\varepsilon_t$ is normal.

\begin{table}[htbp]
\centering
\caption{Critical Values for the Dickey-Fuller $t$-Test}
\label{tab:9.2}
\small
\begin{tabular}{l*{8}{r}}
\toprule
Sample & \multicolumn{8}{c}{Probability that the statistic is less than entry} \\
\cmidrule(l){2-9}
size ($T$) & 0.01 & 0.025 & 0.05 & 0.10 & 0.90 & 0.95 & 0.975 & 0.99 \\
\midrule
\multicolumn{9}{c}{Panel (a): $t$} \\
\midrule
25   & $-2.65$ & $-2.26$ & $-1.95$ & $-1.60$ & 0.92 & 1.33 & 1.70 & 2.15 \\
50   & $-2.62$ & $-2.25$ & $-1.95$ & $-1.61$ & 0.91 & 1.31 & 1.66 & 2.08 \\
100  & $-2.60$ & $-2.24$ & $-1.95$ & $-1.61$ & 0.90 & 1.29 & 1.64 & 2.04 \\
250  & $-2.58$ & $-2.24$ & $-1.95$ & $-1.62$ & 0.89 & 1.28 & 1.63 & 2.02 \\
500  & $-2.58$ & $-2.23$ & $-1.95$ & $-1.62$ & 0.89 & 1.28 & 1.62 & 2.01 \\
$\infty$ & $-2.58$ & $-2.23$ & $-1.95$ & $-1.62$ & 0.89 & 1.28 & 1.62 & 2.01 \\
\midrule
\multicolumn{9}{c}{Panel (b): $t^\mu$} \\
\midrule
25   & $-3.75$ & $-3.33$ & $-2.99$ & $-2.64$ & $-0.37$ & 0.00 & 0.34 & 0.71 \\
50   & $-3.59$ & $-3.23$ & $-2.93$ & $-2.60$ & $-0.41$ & $-0.04$ & 0.28 & 0.66 \\
100  & $-3.50$ & $-3.17$ & $-2.90$ & $-2.59$ & $-0.42$ & $-0.06$ & 0.26 & 0.63 \\
250  & $-3.45$ & $-3.14$ & $-2.88$ & $-2.58$ & $-0.42$ & $-0.07$ & 0.24 & 0.62 \\
500  & $-3.44$ & $-3.13$ & $-2.87$ & $-2.57$ & $-0.44$ & $-0.07$ & 0.24 & 0.61 \\
$\infty$ & $-3.42$ & $-3.12$ & $-2.86$ & $-2.57$ & $-0.44$ & $-0.08$ & 0.23 & 0.60 \\
\midrule
\multicolumn{9}{c}{Panel (c): $t^\tau$} \\
\midrule
25   & $-4.38$ & $-3.95$ & $-3.60$ & $-3.24$ & $-1.14$ & $-0.81$ & $-0.50$ & $-0.15$ \\
50   & $-4.16$ & $-3.80$ & $-3.50$ & $-3.18$ & $-1.19$ & $-0.87$ & $-0.58$ & $-0.24$ \\
100  & $-4.05$ & $-3.73$ & $-3.45$ & $-3.15$ & $-1.22$ & $-0.90$ & $-0.62$ & $-0.28$ \\
250  & $-3.98$ & $-3.69$ & $-3.42$ & $-3.13$ & $-1.23$ & $-0.92$ & $-0.64$ & $-0.31$ \\
500  & $-3.97$ & $-3.67$ & $-3.42$ & $-3.13$ & $-1.24$ & $-0.93$ & $-0.65$ & $-0.32$ \\
$\infty$ & $-3.96$ & $-3.67$ & $-3.41$ & $-3.13$ & $-1.25$ & $-0.94$ & $-0.66$ & $-0.32$ \\
\bottomrule
\multicolumn{9}{l}{\footnotesize\textsc{Source}: Fuller (1976, Table 8.5.2), corrected in Fuller (1996, Table 10.A.2).}
\end{tabular}
\end{table}

\subsection*{Incorporating the Intercept}

A shortcoming of the tests based on the AR(1) equation without intercept is its lack
of invariance to an addition of a constant to the series. If the test is performed on a
series of logarithms (as in Example~9.1 below), then a change in units of measurement of the series results in an addition of a constant to the series, which affects
the value of the test statistic. To make the test statistic invariant to an addition of a
constant, we generalize the model from which the statistic is derived. Consider the
model
\begin{equation}
\label{eq:9.3.10}
y_t = \alpha + z_t, \quad z_t = \rho z_{t-1} + \varepsilon_t, \quad
\{\varepsilon_t\} \text{ independent white noise.}
\end{equation}
The process $\{y_t\}$ obtains when a constant $\alpha$ is added to a process satisfying~\eqref{eq:9.3.1}.
Under the null of $\rho = 1$, $\{z_t\}$ is a driftless random walk. Since $y_t$ can be written as
\[
y_t = \alpha + z_0 + \varepsilon_1 + \varepsilon_2 + \cdots + \varepsilon_t,
\]
$\{y_t\}$ too is a driftless random walk with the initial value of $y_0 \equiv \alpha + z_0$. Under the
alternative of $\rho < 1$, $\{y_t\}$ is stationary AR(1) with mean $\alpha$. Thus, the class of I(0)
processes included in~\eqref{eq:9.3.10} is larger than that in~\eqref{eq:9.3.1}.

By eliminating $\{z_t\}$ from~\eqref{eq:9.3.10}, we obtain an AR(1) equation with intercept
\begin{equation}
\label{eq:9.3.11}
y_t = \alpha^* + \rho y_{t-1} + \varepsilon_t,
\end{equation}
where
\begin{equation}
\label{eq:9.3.12}
\alpha^* = (1 - \rho)\alpha.
\end{equation}
Since $\alpha^* = 0$ when $\rho = 1$, the I(1) null (that the DGP is a driftless random walk) is
the joint hypothesis that $\rho = 1$ \emph{and} $\alpha^* = 0$ in terms of the coefficients of regression~\eqref{eq:9.3.11}. Without the restriction $\alpha^* = 0$, $\{y_t\}$ could be a random walk with drift.
We will develop tests of a random walk with drift in a moment. Until then, we
continue to take the null to be a random walk without drift.

Let $\hat{\rho}^\mu$ here be the OLS estimate of $\rho$ in~\eqref{eq:9.3.11} and $t^\mu$ be the $t$-value for the
null of $\rho = 1$. It should be clear that $\alpha$ would not affect the value of $\hat{\rho}^\mu$ or its OLS
standard error; adding the same constant to $y_t$ for all $t$ merely changes the estimated
intercept. Therefore, the finite-sample as well as large-sample distributions of the
\textbf{DF $\bm{\rho^\mu}$ statistic}, $T \cdot (\hat{\rho}^\mu - 1)$, and the \textbf{DF $\bm{t^\mu}$ statistic}, $t^\mu$, will not depend on $\alpha$
regardless of the value of $\rho$.

As you will be asked to prove in Analytical Exercise~2, $T \cdot (\hat{\rho}^\mu - 1)$ converges
in distribution to a random variable represented as
\begin{equation}
\label{eq:9.3.13}
DF_\rho^\mu \equiv \frac{\tfrac{1}{2}\big([W^\mu(1)]^2 - [W^\mu(0)]^2 - 1\big)}{\int_0^1 [W^\mu(r)]^2\,\mathrm{d}r},
\end{equation}
and $t^\mu$ converges in distribution to
\begin{equation}
\label{eq:9.3.14}
DF_t^\mu \equiv \frac{\tfrac{1}{2}\big([W^\mu(1)]^2 - [W^\mu(0)]^2 - 1\big)}{\sqrt{\int_0^1 [W^\mu(r)]^2\,\mathrm{d}r}}.
\end{equation}
Here, $W^\mu(\cdot)$ is a demeaned standard Wiener process introduced in Section~9.2.
These limiting distributions are free from nuisance parameters such as $\alpha$. The
basic idea of the proof is the following. First, obtain a demeaned series by subtracting the sample mean from $\{y_t\}$ (or equivalently, as the OLS residuals from the
regression of $\{y_t\}$ on a constant). Second, use the Frisch-Waugh Theorem (introduced in Analytical Exercise~4 of Chapter~1) to write $T \cdot (\hat{\rho}^\mu - 1)$ and the $t$-value
in formulas involving the demeaned series. Finally, use Proposition~9.2(c) and~(d).
Summarizing the discussion, we have

\begin{proposition}[Dickey-Fuller tests of a driftless random walk, the case with
intercept]
\label{prop:9.4}
Suppose that $\{y_t\}$ is a driftless random walk (so $\{\Delta y_t\}$ is independent white noise) with $\E(y_0^2) < \infty$. Consider the regression of $y_t$ on $(1, y_{t-1})$ for
$t = 1, 2, \ldots, T$. Then
\begin{align}
&\textit{DF $\rho^\mu$ statistic:}\quad T \cdot (\hat{\rho}^\mu - 1) \;\dto\; DF_\rho^\mu, \notag \\
&\textit{DF $t^\mu$ statistic:}\quad t^\mu \;\dto\; DF_t^\mu, \notag
\end{align}
where $\hat{\rho}^\mu$ is the OLS estimate of the $y_{t-1}$ coefficient, $t^\mu$ is the $t$-value for the
hypothesis that the $y_{t-1}$ coefficient is~$1$, and $DF_\rho^\mu$ and $DF_t^\mu$ are the two random
variables defined in~\eqref{eq:9.3.13} and~\eqref{eq:9.3.14}.
\end{proposition}

Critical values for $DF_\rho^\mu$ are in Table~9.1, Panel~(b), and those for $DF_t^\mu$ are in
Table~9.2, Panel~(b), in rows for $T = \infty$. The distribution is more strongly skewed
to the left than in the case without intercept. The table also has critical values for
finite-sample sizes based on the assumption that $\varepsilon_t$ is normal. The finite-sample
distribution of $\hat{\rho}^\mu$ for $T = 100$ is graphed in Figure~9.3. It shows a greater bias
for $\hat{\rho}^\mu$ than for $\hat{\rho}$ (the OLS estimate without intercept). Unlike in the case without
intercept, the finite-sample critical values do not depend on $y_0$. This is because,
when $\rho = 1$, a change in the initial value $y_0$ only adds the same constant to $y_t$ for
all $t$, so the values of the DF $\rho^\mu$ and $t^\mu$ statistics are invariant to $y_0$.

We started out this subsection by saying that the I(1) null is the joint hypothesis
that $\rho = 1$ \emph{and} $\alpha^* = 0$. Yet the test statistics are for the hypothesis that $\rho = 1$. For
the unit-root tests, this is the practice in the profession; test statistics for the joint
hypothesis are rarely used.

\begin{example}[Are the exchange rates a random walk?]
\label{ex:9.1}
Let $y_t$ be the
log of the spot yen/\$ exchange rate. To make the tests invariant to the choice
of units (e.g., yen/\$ vs.\ yen/cent), we fit the AR(1) model with intercept, so
Proposition~9.4 is applicable. For the weekly exchange rate data used in the
empirical application of Chapter~6 and graphed in Figure~6.3 on page~424,
the AR(1) equation estimated by OLS is
\[
y_t = \underset{(0.435)}{0.162} + \underset{(0.0019285)}{0.9983376}\; y_{t-1},
\quad R^2 = 0.997, \quad SER = 2.824.
\]
Standard errors are in parentheses. Because of the lagged variable $y_{t-1}$, the
actual sample period is from the second week rather than the first week of
1975. So $T = 777$ and
\begin{align*}
T \cdot (\hat{\rho}^\mu - 1) &= 777 \cdot (0.9983376 - 1) = -1.29, \\
t^\mu &= \frac{0.9983376 - 1}{0.0019285} = -0.86.
\end{align*}
Since the alternative hypothesis is that $\rho < 1$, a one-tailed test is called for.
(If you think that the process might be an explosive AR(1) with $\rho > 1$, you
should perform a two-tailed test.) From Table~9.1, Panel~(b), for $DF_\rho^\mu$, the
(asymptotic) critical value that gives 5 percent to the lower tail is $-14.1$.
That is,
\[
\text{Prob}(DF_\rho^\mu < -14.1) = 0.05.
\]
The test statistic of $-1.29$ is well above this critical value, so the null is
easily accepted at the 5 percent significance level. Turning to the DF $t$ test,
from Table~9.2, Panel~(b), the 5 percent critical value is $-2.86$. The $t$-value
of $-0.86$ is greater, so the null hypothesis is easily accepted.
\end{example}

\subsection*{Incorporating Time Trend}

As was mentioned in Section~9.1, most economic time series appear to have linear time trends. To make the DF tests applicable to time series with trend, we
generalize the model once again, this time by adding a linear time trend:
\begin{equation}
\label{eq:9.3.15}
y_t = \alpha + \delta \cdot t + z_t, \quad z_t = \rho z_{t-1} + \varepsilon_t, \quad
\{\varepsilon_t\} \text{ independent white noise.}
\end{equation}
Here, $\delta$ is the slope of the linear trend which may or may not be zero. Under the
null of $\rho = 1$, $\{z_t\}$ is a driftless random walk, and $y_t$ can be written as
\begin{equation}
\label{eq:9.3.16}
y_t = \alpha + \delta \cdot t + z_0 + \varepsilon_1 + \varepsilon_2 + \cdots + \varepsilon_t
= y_0 + \delta \cdot t + (\varepsilon_1 + \cdots + \varepsilon_t)
\end{equation}
with $y_0 \equiv \alpha + z_0$. So $\{y_t\}$ is a random walk with drift if $\delta \neq 0$ and without drift
if $\delta = 0$. Under the alternative that $\rho < 1$, the process, being the sum of a linear
trend and a zero-mean stationary AR(1) process, is trend stationary.

By eliminating $\{z_t\}$ from~\eqref{eq:9.3.15} we obtain
\begin{equation}
\label{eq:9.3.17}
y_t = \alpha^* + \delta^* \cdot t + \rho y_{t-1} + \varepsilon_t,
\end{equation}
where
\begin{equation}
\label{eq:9.3.18}
\alpha^* = (1 - \rho)\alpha + \rho\delta, \quad \delta^* = (1 - \rho)\delta.
\end{equation}
Since $\delta^* = 0$ when $\rho = 1$, the I(1) null (that the DGP is a random walk with
or without drift) is the joint hypothesis that $\rho = 1$ \emph{and} $\delta^* = 0$ in terms of the
regression coefficients. In practice, however, unit-root tests usually focus on the
single restriction $\rho = 1$.

As in the AR(1) model with intercept, the statement of Proposition~9.3 can be
readily modified to the AR(1) model with trend. Let $\hat{\rho}^\tau$ be the OLS estimate of the
$\rho$ in~\eqref{eq:9.3.17} and $t^\tau$ be the $t$-value for the null of $\rho = 1$. As you will be asked to
prove in Analytical Exercise~3, $T \cdot (\hat{\rho}^\tau - 1)$ converges in distribution to
\begin{equation}
\label{eq:9.3.19}
DF_\rho^\tau \equiv \frac{\tfrac{1}{2}\big([W^\tau(1)]^2 - [W^\tau(0)]^2 - 1\big)}{\int_0^1 [W^\tau(r)]^2\,\mathrm{d}r},
\end{equation}
and $t^\tau$ converges in distribution to
\begin{equation}
\label{eq:9.3.20}
DF_t^\tau \equiv \frac{\tfrac{1}{2}\big([W^\tau(1)]^2 - [W^\tau(0)]^2 - 1\big)}{\sqrt{\int_0^1 [W^\tau(r)]^2\,\mathrm{d}r}}.
\end{equation}
Here, $W^\tau(\cdot)$ is a detrended standard Wiener process introduced in Section~9.2. The
basic idea of the proof is to use the Frisch-Waugh Theorem to write $T \cdot (\hat{\rho}^\tau - 1)$ and
$t^\tau$ in formulas involving detrended series and then apply Proposition~9.2(e) and~(f).
Summarizing the discussion, we have

\begin{proposition}[Dickey-Fuller tests of a random walk with or without drift]
\label{prop:9.5}
Suppose that $\{y_t\}$ is a random walk with or without drift with $\E(y_0^2) < \infty$. Consider the regression of $y_t$ on $(1, t, y_{t-1})$ for $t = 1, 2, \ldots, T$. Then
\begin{align}
&\textit{DF $\rho^\tau$ statistic:}\quad T \cdot (\hat{\rho}^\tau - 1) \;\dto\; DF_\rho^\tau,
\label{eq:9.3.21} \\
&\textit{DF $t^\tau$ statistic:}\quad t^\tau \;\dto\; DF_t^\tau,
\label{eq:9.3.22}
\end{align}
where $\hat{\rho}^\tau$ is the OLS estimate of the $y_{t-1}$ coefficient, $t^\tau$ is the $t$-value for the hypothesis that the $y_{t-1}$ coefficient is~$1$, and $DF_\rho^\tau$ and $DF_t^\tau$ are the two random variables
defined in~\eqref{eq:9.3.19} and~\eqref{eq:9.3.20}.
\end{proposition}

Two comments are in order about Proposition~9.5.

\begin{itemize}
\item (Numerical invariance) \enspace The test statistics $\hat{\rho}^\tau$ and $t^\tau$ are invariant to $(\alpha, \delta)$ regardless of the value of $\rho$; adding $\alpha + \delta \cdot t$ to the series $\{y_t\}$ merely changes the OLS
estimates of $\alpha^*$ and $\delta^*$ but not $\hat{\rho}^\tau$ or its $t$-value. Therefore, the finite-sample
as well as large-sample distributions of $T \cdot (\hat{\rho}^\tau - 1)$ and $t^\tau$ will not depend on
$(\alpha, \delta)$, regardless of the value of $\rho$.

\item (Should the regression include time?) \enspace Since the test statistics are invariant to
the value $\delta$ in the DGP, Proposition~9.5 is applicable even when $\delta = 0$. That is,
the $\hat{\rho}^\tau$ and $t^\tau$ statistics can be used to test the null of a \emph{driftless} random walk.
However, if you are confident that the null process has no trend, then the DF $\rho^\mu$
and $t^\mu$ tests of Proposition~9.4 should be used, with the regression not including
time, because the finite-sample power against stationary alternatives is generally
greater with $\hat{\rho}^\mu$ and $t^\mu$ than with $\hat{\rho}^\tau$ and $t^\tau$ (you will be asked to verify this in
Monte Carlo Exercise~1). On the other hand, if you think that the process might
have a trend, then you should use $\hat{\rho}^\tau$ and $t^\tau$, with time in the regression and use critical values from Table~9.1, Panel~(c), and Table~9.2, Panel~(c). If you
fail to include time in the regression and use critical values from Table~9.1, Panel~(b), and Table~9.2, Panel~(b), when indeed the DGP has a trend, then the test will
not be correctly sized in large samples (i.e., the probability of rejecting the null
when the null is true will not approach the prespecified significance level (the
nominal size) as the sample size goes to infinity). This is simply because the
limiting distributions of $\hat{\rho}^\mu$ or $t^\mu$ when the DGP has a trend are not $DF_\rho^\mu$ or
$DF_t^\mu$.
\end{itemize}

Critical values for $DF_\rho^\tau$ are in Table~9.1, Panel~(c), and those for $DF_t^\tau$ are in Table~9.2, Panel~(c). The table also has critical values for finite sample sizes based on the
assumption that $\varepsilon_t$ is normal (the initial value $y_0$ as well as the trend parameters $\alpha$
and $\delta$ need not be specified to calculate the finite-sample distributions under the null of $\rho = 1$). For
both the DF $\rho^\tau$ and $t^\tau$ statistics, the distribution is even more strongly skewed to
the left than that for the case with intercept. Now for all sample sizes displayed in the
tables, the 99 percent critical value is negative, meaning that more than 99 percent
of the time $\hat{\rho}^\tau$ is less than unity when in fact the true value of $\rho$ is unity! This
is illustrated by the chained curve in Figure~9.3, which graphs the finite-sample
distribution of $\hat{\rho}^\tau$ for $T = 100$.

\bigskip
\noindent\textsc{questions for review} \hrulefill

\begin{enumerate}
\item (Superconsistency) \enspace Let $\hat{\rho}$ be the OLS coefficient estimate as in~\eqref{eq:9.3.2}. Verify
that $T^{1-\eta} \cdot (\hat{\rho} - 1) \pto 0$ for $0 < \eta < 1$ if $\rho = 1$.

\item (Limiting distribution under general driftless I(1)) \enspace In the AR(1) model~\eqref{eq:9.3.1},
drop the assumption that the error term is independent white noise by replacing
$\varepsilon_t$ by $u_t$ and assuming that $\{u_t\}$ is a general zero-mean I(0) process satisfying
\eqref{eq:9.2.1}--\eqref{eq:9.2.3b}. So the process under the null of $\rho = 1$ is a general driftless
I(1) process. Show that $T \cdot (\hat{\rho} - 1)$ converges in distribution to
\[
\frac{\tfrac{1}{2}\big(W(1)^2 - \tfrac{\gamma_0}{\lambda^2}\big)}{\int_0^1 W(r)^2\,\mathrm{d}r}
\]
where $\lambda^2 \equiv \text{long-run variance of } \{\Delta y_t\}$ and $\gamma_0 \equiv \Var(\Delta y_t)$.

\item (Consistency of $s^2$) \enspace The usual OLS standard error of regression for the AR(1)
equation, $s$, is defined in~\eqref{eq:9.3.8}. If $|\rho| < 1$, we know from Section~6.4 that
$s^2$ is consistent for $\Var(\varepsilon_t)$. Show that, under the null of $\rho = 1$, $s^2$ remains
consistent for $\gamma_0$ ($\equiv \Var(\Delta y_t)$) as well. \textbf{Hint:} Rewrite $s^2$ as
\begin{align*}
s^2 &\equiv \frac{1}{T-1}\sum_{t=1}^{T} (\Delta y_t - (\hat{\rho}-1) y_{t-1})^2 \\
&= \frac{1}{T-1}\sum_{t=1}^{T}(\Delta y_t)^2
- \frac{2}{T-1} \cdot [T \cdot (\hat{\rho}-1)] \cdot \frac{1}{T}\sum_{t=1}^{T}\Delta y_t\, y_{t-1} \\
&\quad + \frac{1}{T-1} \cdot [T \cdot (\hat{\rho}-1)]^2 \cdot \frac{1}{T^2}\sum_{t=1}^{T}(y_{t-1})^2.
\end{align*}
Also, $\E[(\Delta y_t)^2] = \gamma_0$. Use the result we have used repeatedly: if $x_T \pto 0$ and
$y_T \dto y$, then $x_T \cdot y_T \pto 0$.

\item (Two tests for a random walk) \enspace Consider two regressions. One is to regress
$\Delta y_t$ on $\Delta y_{t-1}$, and the other is to regress $\Delta y_t$ on $y_{t-1}$. To test the null that $\{y_t\}$
is a driftless random walk, which distribution should be used, the $t$-value from
the first regression or the $t$-value from the second regression?

\item (Consistency of DF tests) \enspace Recall that a test is said to be \textbf{consistent} against a
set of alternatives if the probability of rejecting a false null hypothesis when
the true DGP is any of the alternatives approaches unity as the sample size
increases. Let $T \cdot (\hat{\rho} - 1)$ be as defined in~\eqref{eq:9.3.3}.

\begin{enumerate}[label=(\alph*)]
\item Suppose that $\{y_t\}$ is a zero-mean I(0) process and that its first-order autocorrelation coefficient is less than~1 (so $\E(y_t^2) > \E(y_t\, y_{t-1})$). Show that
\[
\plim_{T \to \infty} T \cdot (\hat{\rho} - 1) = -\infty.
\]
Thus, the DF $\rho$ test is consistent against general I(0) alternatives. \textbf{Hint:}
Show that, if $\{y_t\}$ is I(0),
\[
\plim_{T \to \infty} \hat{\rho} = \frac{\E(y_t\, y_{t-1})}{\E(y_{t-1})^2}.
\]

\item Show that the DF $t$ test is consistent against general zero-mean I(0) alternatives. \textbf{Hint:} $s$ converges in probability to some positive number when $\{y_t\}$
is zero-mean I(0).
\end{enumerate}
\end{enumerate}
